\introsection{Цели и задачи исследования}

В ходе планирования реализации концепта интеллектуального модуля интерактивного анализа данных в реляционных базах данных была выдвинута цель: \textbf{«Разработать программное решение, которое устранит проблему обращения к большим данным, пользователям без должного технического образования»}. В ходе реализации был разработан NLIDB интерфейс, который закроет проблемы пользователей.

Исходя из цели, были выбраны следующие задачи:

\begin{enumerate}
    \item Провести оценку современных методов трансляции естественного языка в язык запросов;

    \item Спроектировать RAG-архитектуру, для работы с таблицами в схеме БД, с выявлением метаинформации, связи таблиц;

    \item Разработка механизма самокоррекции ИИ-агента, который позволит улучшить качество готового ответа;

    \item Провести оценку качества генерации, при добавлении в проект механизма самокоррекции;

    \item Реализовать микросервисную архитектуру системы.

\end{enumerate}
